\subsection{A cyclotomic criterion and the remaining symmetric covers}

\begin{statement}{A.38}{Proved}{The nonstable symmetric genus-three fixed groups at every prime}
For \(a=6,-6\), put
\[
 C_a:v^2=w^7+aw^4+w,\qquad
 X_a=C_a\setminus\{P_0,P_\infty\}.
\]
For every prime \(\ell\), every original packet of Part A,
and every permitted open subgroup \(\Delta_0\) preserving
the cover and the selected cusp lift \(P\), one has
\[
 \left(\pi_1(X_{a,\overline K})^{(\ell)}\right)^{s_b(\Delta_0)}
                =I_{P,\ell}.                         \tag{A.38.1}
\]
At odd primes the argument also applies to every cusp
packet over a finite number field of definition, with
that field's own derived Galois subgroup. It includes
the packet with zero first Johnson coefficient.
No nonzero primitive Collino regulator is required.

More generally, let \(X/k\) be a hyperbolic affine curve
over a number field with at least two geometric cusps,
let \(s\) be a packet at a rational cusp, and write
\(F=\pi_1(X_{\bar k})^{(\ell)}\),
\(H=F^{s(\overline{[G_k,G_k]})}\).
Give its rational unipotent completion the weight
filtration in which proper homology has weight one
and boundary generators have weight two. Suppose the
Lie algebra of the unipotent closure of \(H\) has
zero odd graded pieces, and \(G_k\) acts on its
weight \(2m\) piece by the scalar \(\chi_\ell^m\).
Then \(H\) is exactly the selected cusp inertia.
This criterion concerns the closure of an actual
subgroup; it does not assert that a rational fixed
Lie element lifts to \(F\).
\end{statement}

\begin{proof}
\paragraph{The arithmetic criterion.}
We first prove the general assertion, keeping the
given field and the given actual subgroup \(H\).
The actual packet makes \(G_k\) act on \(F\), and
\(H\) is stable under this action. On \(H\) the
action factors through \(G_k^{\rm ab}\), since the
closed commutator subgroup fixes every element of
\(H\). This remains true on every algebraic unipotent
closure of its image.

Let \(\mathfrak h\) denote the inverse limit of the
Lie algebras of those closures in the finite weight
quotients of the rational unipotent completion of
\(F\). Images of homomorphisms of unipotent algebraic
groups in characteristic zero are closed. Thus the
transition maps between these closures are surjective,
and their Lie algebras are the finite quotients of
\(\mathfrak h\) with its induced weight filtration.
All these quotients are finite dimensional.

Choose a finite place of residue cardinality \(q\ne1\)
and residue characteristic different from \(\ell\),
and choose any arithmetic Frobenius lift \(\varphi\).
We use only \(\chi_\ell(\varphi)=q\). On a finite
filtered quotient of \(\mathfrak h\), the action of
\(\varphi\) on weight \(2m\) is the scalar \(q^m\).
These scalars are distinct for distinct \(m\).
A filtered linear operator which is scalar on every
graded piece with distinct such scalars is semisimple:
its minimal polynomial divides the product of the
distinct factors \(T-q^m\), as follows successively
through the filtration. Its eigenspace for \(q^m\)
maps isomorphically to the weight \(2m\) graded piece.

Now let \(\gamma\in\ker\chi_\ell\). It commutes
with \(\varphi\) on \(\mathfrak h\), because the
action there is abelian. It induces the identity on
every associated graded piece. It therefore fixes
each \(\varphi\)-eigenspace pointwise, and acts as
the identity on every finite quotient of
\(\mathfrak h\). Exponentiation gives the same
conclusion on its unipotent groups.

The map from the finite-rank free pro-\(\ell\)
group \(F\) to its rational unipotent completion is
injective. Indeed its integral Magnus expansion is
injective, its ordinary lower-central graded pieces
are torsion free, and the weighted filtration is
separated: assigning only weights one and two makes
it comparable with the ordinary degree filtration.
We conclude for the actual group that
\[
                     H\subset F^{s(\ker\chi_\ell)}.
                                                            \tag{A.38.2}
\]
This uses rational completions only to establish an
identity on existing integral elements.

We next prove that every element of the right-hand
side of (A.38.2) is peripheral. Suppose \(h\) were
not peripheral. The purely group-theoretic detection
lemma of \cite[Lemma 2.1.4, printed pp.96--97]{Nakamura1994},
applied to the full class of finite \(\ell\)-groups,
gives an open subgroup \(U\subset F\) containing
\(h\) in which its image in proper homology modulo
\(\ell\) is nonzero. More explicitly,
\[
 h\notin\overline{[U,U]U^\ell
                  \langle Z\cap U\rangle},          \tag{A.38.3}
\]
where \(Z\) is the union of all conjugates of all
cusp inertia groups. Killing \(Z\cap U\) is exactly
proper compactification of the corresponding cover.
The proper integral homology is a free
\(\mathbf Z_\ell\)-module, so (A.38.3) also gives
a nonzero rational homology vector.

There are finitely many subgroups of the given index
in the finitely generated group \(F\). After a finite
extension \(k'/k\), its actual packet action preserves
\(U\). The group \(U\rtimes s(G_{k'})\) is then
an actual arithmetic finite-cover subgroup. The
corresponding geometric cover has \(\ell\)-group
Galois closure; the core argument of A.22 identifies
its full geometric pro-\(\ell\) group with \(U\).
Let \(J\) be the Jacobian of its smooth proper
compactification. The vector just found belongs to
\(V_\ell J\) and, by (A.38.2), is fixed by
\(G_{k'(\mu_{\ell^\infty})}\).

This contradicts \cite[Theorem 3, printed p.315,
proof pp.317--318]{A_RibetCyclotomic1981}.
That theorem gives
\[
       (V_\ell J)^{G_{k'(\mu_\infty)}}=0             \tag{A.38.4}
\]
for every abelian variety over every number field
and every prime. The subgroup fixing all roots of
unity is contained in the subgroup fixing the
\(\ell\)-power roots, so (A.38.4) gives the required
vanishing as well. All hypotheses concern the
Jacobian of this actual finite cover; no claim is
made about homology of an infinite cover.

Thus every nonidentity element of \(H\) is peripheral.
Since \(X\) has at least two geometric cusps, every
cusp has a nonzero primitive class in ordinary
geometric pro-\(\ell\) abelianization. Every nonzero
power of an inertia generator therefore has nonzero
ordinary homology. Every commutator of two elements
of \(H\), on the other hand, has zero ordinary
homology. It must be the identity. Hence \(H\) is
abelian. It contains the selected inertia \(I\),
because the cyclotomic character is trivial on
the closed commutator subgroup. The centralizer
of a nontrivial element of this primitive free
factor is \(I\) itself, by
\cite[Lemma 2.1.2]{Nakamura1994}. Consequently
\(H=I\), proving the criterion.

\paragraph{The weighted monodromy for the two symmetric curves.}
We apply this criterion to the curves in the statement
and odd \(\ell\). Choose a finite field of definition
\(k\) so that
\(\Gamma=\overline{[G_k,G_k]}\subset\Delta_0\).
All marked and deck data of A.33 are defined there.
It is enough to compute the \(\Gamma\)-fixed group:
it contains the desired \(\Delta_0\)-fixed group,
and \(\Delta_0\) already fixes the selected inertia.

A.33 establishes the actual pointed reduction at
two: a rational centre bearing the two cusps, with
three elliptic tails. Let \(t\) denote the logarithm
of the selected primitive cusp generator, and let
\(a_i,b_i\) be the six handle generators attached
to these tails. In the weight-graded rational Lie
algebra put
\[
 \operatorname{wt}(a_i)=\operatorname{wt}(b_i)=1,
 \quad \operatorname{wt}(t)=2,
 \quad \omega_i=[a_i,b_i],\quad
 \Omega=\omega_1+\omega_2+\omega_3.                  \tag{A.38.5}
\]
It is the free weighted Lie algebra on the handle
space and \(t\). This filtration is distinct from
the ordinary lower-central filtration used in A.33.

For an ordinary base point on the rational centre,
the first nonzero weighted monodromy derivation is
\(D_0v_i=m_i[\omega_i,v_i]\), \(D_0t=0\), with
all \(m_i\ne0\). This is the actual based edge
formula of A.33 and Betts--Dogra. The edges enclose
only their indicated genus-one components. Their
boundary logarithms have first weight two term
\(\omega_i\); higher terms cannot change this
degree-two derivation.

Compare the ordinary point with the selected tangent
through its punctured formal residue disc, as in A.36.
The comparison is inner by the selected cusp group.
Equivalently, at weight two an inner discrepancy
fixing \(t\) has conjugator in the degree-two
centralizer of \(t\), which is exactly
\(\mathbf Q_\ell t\). There is no weight-one
inner discrepancy fixing \(t\). Thus for every
packet the genuine derivation of an actual local
Frobenius--tame-inertia commutator in \(\Gamma\)
is a nonzero common multiple of
\[
 \begin{split}
 D_\lambda(t)&=0,\\
 D_\lambda(v_i)&=[m_i\omega_i+\lambda t,v_i]
       \quad(v_i=a_i,b_i),
 \qquad\lambda\in\mathbf Q_\ell.
 \end{split}
                                                        \tag{A.38.6}
\]
Both summands now have weight shift two. In
particular the zero first ordinary Johnson case
does not require a separate nonvanishing argument.

\paragraph{The unbounded weighted kernel.}
Inside the free associative enveloping algebra set
\(B=\mathbf Q_\ell\langle t,\omega_1,\omega_2,
\omega_3\rangle\). These are freely independent
associative generators: introduce symbols \(u_i\)
and orient
\[
                  a_i b_i\longmapsto b_i a_i+u_i.
\]
The leading words have no overlaps, and the normal
words give an injective copy of the free algebra
on \(t,u_1,u_2,u_3\). This is the same normal-form
construction as A.27, with the additional letter
\(t\).

Group normal words by their ordered sequence of
handle letters, and place the words in \(B\)
in the intervening slots. On \(B\), (A.38.6)
is \(\lambda\operatorname{ad}t\). For a nonempty
handle sequence beginning in the \(i\)-th block,
the operator increasing its first slot's length is
\[
                       \lambda L_t+m_iR_{u_i}.       \tag{A.38.7}
\]
It is injective. For \(\lambda=0\) this is right
multiplication by a free letter; for
\(\lambda\ne0\), the greatest number of occurrences
of \(t\) in a finite input is increased only by the
left summand. All other parts of the derivation
preserve that first-slot length. A component of
greatest first-slot length in a finite polynomial
therefore cannot lie in the kernel. There is no
kernel outside \(B\).

Within \(B\) the kernel is all of \(B\) when
\(\lambda=0\), and \(\mathbf Q_\ell[t]\) otherwise.
Intersecting with primitive elements gives, in
particular,
\[
 \ker D_\lambda\subset
               \mathbb L(t,\omega_1,\omega_2,\omega_3).
                                                            \tag{A.38.8}
\]
This is a proof in every weight. It applies to
the associated graded of the unipotent closure
of actual \(\Gamma\)-fixed elements, because
the commutator defining (A.38.6) fixes that closure
pointwise.

The non-CM invariant intersection of A.33 is unchanged
in the weighted grading. We recall the exact linear
map to make this point explicit. Put
\(A=\sum a_i\), \(B_1=\sum b_i\), and
\(\delta a_i=a_i-A/3\). The Zariski closure of
\(\Gamma\) on the two copies of the non-CM elliptic
factor contains the symplectic transformation
\[
 g(a_i)=a_i,\qquad g(b_i)=b_i+\delta a_i,\qquad g(t)=t.
\]
For \(U=\langle\omega_1,\omega_2,\omega_3\rangle\)
one computes
\[
 g(\omega_i)=\omega_i+\tfrac13[A,\delta a_i],
 \qquad U\cap gU=\mathbf Q_\ell\Omega.               \tag{A.38.9}
\]
The second identity follows by comparing the
\(a_i\wedge a_j\) components: the sum-zero vector
\(\sum d_i\delta a_i\) must commute with \(A\),
and hence vanish, so all \(d_i\) are equal.
In associative words, split at each occurrence
of \(t\), and pair consecutive handle letters
in each intervening block. Tensor-space intersection
in each pair replaces \(U\) and \(gU\) by their
intersection. Taking primitives yields
\[
 \mathbb L(t,U)\cap\mathbb L(t,gU)
                       =\mathbb L(t,\Omega).         \tag{A.38.10}
\]
Both \(t\) and \(\Omega\) have weight two.

The full arithmetic action on these two generators
is \(\chi_\ell\): for \(t\) this is actual cusp
inertia, and for \(\Omega\) it is the symplectic
similitude character on proper homology. It follows
from (A.38.8)--(A.38.10) that every weight \(2m\)
graded piece of the fixed subgroup's unipotent
closure has scalar character \(\chi_\ell^m\),
and every odd piece vanishes. The criterion proved
above applies and gives exactly \(I_{P,\ell}\).

At \(\ell=2\), A.33 already gives the assertion
by the actual unramified double-cover comparison.
All steps at odd primes are over \(\mathbf Q_\ell\)
only for detecting identities on existing elements;
the final cusp group and equality are integral.
This includes \(\ell=3\), despite the harmless
rational denominators in (A.38.9).

The executed certificate \path{code/a_cyclotomic_fixed.py}
checks the first-slot operators over \(\mathbf Q\)
for the four boundary letters through slot length
four, and checks a filtered Frobenius matrix with
four two-dimensional scalar graded pieces and
nonzero extension blocks. Its commutation equations
have no nonzero strictly filtration-increasing
solution. These finite checks support the preceding
unbounded proof; they are not used to realize
rational Lie classes as actual group elements.
\end{proof}
